\documentclass[12pt, a4paper]{article}

% --- Настройки для русского языка под стандартный pdfLaTeX ---
\usepackage[T2A]{fontenc}
\usepackage[utf8]{inputenc}
\usepackage[english, russian]{babel}

% --- Математические пакеты ---
\usepackage{amsmath, amsfonts, amssymb}

% --- Настройка полей ---
\usepackage[left=2cm, right=2cm, top=2cm, bottom=2cm]{geometry}

\begin{document}

\section*{Ответы на вопросы к коллоквиуму (МОПС)}

\begin{enumerate}
    \item \textbf{Записать определение вероятности.} \\
    Вероятность — это количественная мера (степень) возможности наступления некоторого случайного события. В классическом определении $P(A) = m/n$, где $m$ — число благоприятствующих исходов, а $n$ — общее число равновозможных исходов.

    \item \textbf{Записать определение плотности вероятности.} \\
    Плотность вероятности $p(x)$ (для непрерывной случайной величины) — это производная от функции распределения вероятностей $F(x)$, то есть $p(x) = \frac{dF(x)}{dx}$.

    \item \textbf{Определение совместной плотности вероятности.} \\
    Для двух случайных величин $X$ и $Y$ совместная плотность вероятности — это вторая смешанная частная производная от их совместной функции распределения: $p(x, y) = \frac{\partial^2 F(x,y)}{\partial x \partial y}$.

    \item \textbf{Определение условной плотности вероятности.} \\
    Плотность распределения одной случайной величины при условии, что другая приняла заданное значение: $p(x|y) = \frac{p(x,y)}{p(y)}$.

    \item \textbf{Записать формулу Байеса.} \\
    $p(x|y) = \frac{p(y|x)p(x)}{p(y)}$

    \item \textbf{Записать определение плотности вероятности многомерной случайной величины.} \\
    Это $n$-я смешанная частная производная совместной функции распределения по всем координатам: $p(x_1, \dots, x_n) = \frac{\partial^n F(x_1, \dots, x_n)}{\partial x_1 \dots \partial x_n}$.

    \item \textbf{Определение математического ожидания многомерной случайной величины.} \\
    Это детерминированный вектор, компонентами которого являются математические ожидания соответствующих компонент случайного вектора.

    \item \textbf{$x$ - случайный вектор-столбец размерности $n$. Какова размерность $M[x]$?} \\
    Размерность вектора $M[x]$ равна $n \times 1$ (вектор-столбец).

    \item \textbf{$x$ - случайный вектор-столбец размерности $n$. Какова размерность $p(x)$?} \\
    Размерность $1 \times 1$ (скаляр).

    \item \textbf{$Y$ - матрица случайных чисел размером $2048 \times 4$. Какова размерность $M[Y]$?} \\
    Размерность останется $2048 \times 4$.

    \item \textbf{Определение дисперсии многомерной случайной величины.} \\
    Аналогом является ковариационная матрица: $D = M[(x - M[x])(x - M[x])^T]$.

    \item \textbf{Назвать свойства матрицы дисперсий.} \\
    1) Симметричность ($D = D^T$). 2) Неотрицательная определенность. 3) На главной диагонали — дисперсии компонент, вне диагонали — их взаимные ковариации.

    \item \textbf{$x$ - случайный вектор-столбец размерности $n$. Какова размерность его матрицы дисперсий?} \\
    Размерность $n \times n$ (квадратная матрица).

    \item \textbf{$x$ - вектор-столбец размерности $n$. Чему равна размерность величины $A=x \cdot x^T$?} \\
    Вектор-столбец ($n \times 1$) умножается на вектор-строку ($1 \times n$). Результат — матрица $n \times n$.

    \item \textbf{$x$ - вектор-столбец размерности $n$. Чему равна размерность величины $A=x^T \cdot x$?} \\
    Вектор-строка ($1 \times n$) умножается на вектор-столбец ($n \times 1$). Результат — скаляр размерности $1 \times 1$.

    \item \textbf{$a=\begin{bmatrix}1 & 2 & 3\end{bmatrix}$, $b=\begin{bmatrix}4 \\ 5\end{bmatrix}$}
    \begin{itemize}
        \item \textbf{Чему равно ba?} Умножается вектор ($2 \times 1$) на вектор ($1 \times 3$). Получится матрица ($2 \times 3$): 
        $ba = \begin{bmatrix} 4 \\ 5 \end{bmatrix} \begin{bmatrix} 1 & 2 & 3 \end{bmatrix} = \begin{bmatrix} 4 & 8 & 12 \\ 5 & 10 & 15 \end{bmatrix}$
        \item \textbf{Имеет ли смысл выражение ab?} Не имеет. Внутренние размерности (3 и 2) не совпадают.
        \item \textbf{Имеет ли смысл выражение a + b?} Не имеет. Сложение матриц с разными размерностями невозможно.
    \end{itemize}

    \item \textbf{Матрица $A$ имеет размерность $3 \times 3$. Вектор $b$ имеет размерность $3 \times 1$. Какую будет иметь размерность величина:}
    \begin{itemize}
        \item а) $A \cdot b$ --- Размерность $3 \times 1$ (вектор-столбец).
        \item б) $b \cdot A$ --- Не имеет смысла.
        \item в) $b^T \cdot A \cdot b$ --- Размерность $1 \times 1$ (скаляр/квадратичная форма).
    \end{itemize}

    \item \textbf{Матрица $A$ имеет размерность $2 \times 4$. Матрица $B$ имеет размерность $4 \times 7$. Какую размерность будет иметь матрица $A \cdot B$?} \\
    Размерность будет $2 \times 7$.

    \item \textbf{Дана векторная функция $f$ (размерности $m$) векторного аргумента $x$ (размерности $n$). Чему по определению равна производная $\frac{\partial f(x)}{\partial x}$?} \\
    Это матрица Якоби размерностью $m \times n$. Элемент матрицы равен $\frac{\partial f_i}{\partial x_j}$.

    \item \textbf{Векторная функция $f$ имеет размерность 10, её векторный аргумент $x$ имеет размерность 3. Какую размерность имеет производная?} \\
    Матрица Якоби будет иметь размерность $10 \times 3$.

    \item \textbf{$x = \begin{bmatrix}x \\ y \\ z\end{bmatrix}$, $f(x) = \begin{bmatrix} A \cdot \sin(2x-y) \\ A \cdot \cos(3z-y) \end{bmatrix}$. Найти $\frac{\partial f(x)}{\partial x}$} \\
    Матрица Якоби (размерность $2 \times 3$):
    \[
    \begin{bmatrix} 
    2A\cos(2x-y) & -A\cos(2x-y) & 0 \\. 
    0 & A\sin(3z-y) & -3A\sin(3z-y) 
    \end{bmatrix}
    \]

    \item \textbf{$n$ - случайный вектор ($M[n]=m_n$, $D_n$), $a$ - известное число.} \\
    $M[an] = a \cdot m_n$, \quad $D[an] = a^2 \cdot D_n$

    \item \textbf{$n$ - случайный вектор... $S$ - известный вектор. Чему равняется матожидание и дисперсия величины $(S+n)$?} \\
    $M[S+n] = S + m_n$, \quad $D[S+n] = D_n$

    \item \textbf{Дана выборка $x:|-18 \quad 4 \quad 2|$. Найти выборочное среднее.} \\
    $\bar{x} = \frac{-18 + 4 + 2}{3} = -4$

    \item \textbf{Дана выборка $x:|-18 \quad 4 \quad 2|$. Найти выборочную дисперсию.} \\
    Несмещенная выборочная дисперсия:
    $S^2 = \frac{(-18 - (-4))^2 + (4 - (-4))^2 + (2 - (-4))^2}{3 - 1} = \frac{196 + 64 + 36}{2} = 148$

    \item \textbf{Что означает запись $\int_x f(x)dx$ если $f$ - векторная функция векторного аргумента $x$?} \\
    Это многомерный интеграл. Операция интегрирования применяется к каждому элементу вектор-функции $f(x)$ отдельно по всему пространству параметров вектора $x$. Результатом будет вектор-столбец.

\end{enumerate}

\end{document}